\documentclass{article}
\usepackage{amsmath,amssymb,amsthm}
\usepackage{graphicx}
\usepackage{hyperref}

\title{Theory of Incidence -- A Fourth Foundation Beyond Set, Category, and Type}
\author{Junkawasaki}
\date{\today}

\begin{document}

\maketitle

\begin{abstract}
Classical mathematical foundations—set theory, category theory, and type theory—provide distinct but compatible paradigms for representing mathematical and computational structure. Yet all share a fundamental limitation: they presuppose a separation between objects and relations. Sets collect elements, categories connect objects, and types classify terms, but none take relation itself as the primitive entity.

We propose the Theory of Incidence as a fourth foundational system in which incidences—relations that may themselves relate to other incidences—are the sole primitive constituents of mathematical structure. Each incidence $i \in I$ is equipped with a boundary operator $\partial: I \to MSet(I \times R \times \Sigma \times N)$, assigning to $i$ a finite or guardedly infinite multiset of endpoint incidences, roles, orientations, and multiplicities. Composition of structures is governed by a gluing operator, subject to axioms (A1–A17) ensuring finiteness, type-consistency, compositional closure, and observational equivalence.

Within this framework,
\begin{itemize}
\item Sets correspond to zero-ary incidences,
\item Categories arise as incidence systems closed under gluing, and
\item Types emerge as inductive or coinductive families over incidence configurations.
\end{itemize}

The theory admits multiple semantic realizations: a set-theoretic model (conservative over ZF), an adhesive-category model (gluing semantics), and a homotopy-type-theoretic model (via coinductive types). A linear-semantic layer—defined through boundary matrices and Laplacians—connects logical inference to spectral computation, linking proof theory with numerical analysis.

The Theory of Incidence thus provides a unifying relational substrate for Set, Cat, and Type: a formalism where relations are first-class citizens, recursion and coinduction are native, and evolving or self-referential structures are expressible without additional meta-level machinery. This approach suggests a direction toward a unified foundation of mathematics and computation based purely on relational composition and equivalence.
\end{abstract}

\section{Introduction}
Classical foundations separate objects from relations. Incidence Theory unifies them by making relations primitive.

\section{Formal Definitions}
\begin{itemize}
\item Incidence set $I$
\item Boundary operator $\partial: I \to MSet(I \times R \times \Sigma \times N)$
\end{itemize}

\section{Axioms (A1–A17)}
\begin{table}[h]
\centering
\caption{Axioms of Incidence Theory (A1–A17)}
\label{tab:axioms}
\begin{tabular}{|c|l|l|}
\hline
Axiom & Description & Lean Implementation \\
\hline
A1 & Finite Endpoints & \texttt{Incidence.boundary\_finite} \\
A2 & Type Consistency & \texttt{Incidence.type\_consistent} \\
A3 & Sign Rules & \texttt{Incidence.boundary\_signs} \\
A4 & Multiplicities & \texttt{Incidence.boundary\_multiplicities} \\
A5 & Well-founded Mode & \texttt{Incidence.well\_founded} \\
A6 & Gluing Existence & \texttt{Incidence.gluing} \\
A7 & Unit Laws & \texttt{Incidence.gluing\_unit\_left/right} \\
A8 & Associativity & \texttt{Incidence.gluing\_associative} \\
A9 & Type Preservation & \texttt{Incidence.gluing\_type\_preserves} \\
A10 & Guard Preservation & \texttt{Incidence.guard\_preservation} \\
A11 & Observational Equivalence & \texttt{approx}, \texttt{approx\_refl/symm/trans} \\
A12 & Congruence & \texttt{Incidence.congruence\_gluing} \\
A13 & Normalization & \texttt{normalization\_property} \\
A14 & Functoriality & \texttt{functoriality} \\
A15 & Colimits & \texttt{colimits\_exist} \\
A16 & Boundary Matrix & \texttt{boundary\_matrix\_well\_defined} \\
A17 & Laplacian & \texttt{laplacian\_well\_defined} \\
\hline
\end{tabular}
\end{table}

\section{Relations to Set, Category, and Type Theories}

We establish embeddings of set theory, category theory, and type theory into the Theory of Incidence.

\begin{theorem}[Set Embedding]
\label{thm:set-embedding}
Sets can be represented as incidences containing nullary incidences as elements, with extensionality preserved through observational equivalence.
\end{theorem}

\begin{proof}
A set $S$ is represented by an incidence $i_S$ where:
\begin{itemize}
\item The boundary of $i_S$ contains exactly the nullary incidences representing elements of $S$
\item Membership $x \in S$ corresponds to $x$ appearing in the boundary of $i_S$
\item Extensionality: Two sets are equal iff they have the same elements
\end{itemize}

This embedding is formalized in Lean as \texttt{represents\_set} and \texttt{set\_embedding\_extensionality} in \texttt{IncidenceTheory/Basic.lean}.
\end{proof}

\begin{theorem}[Category Embedding]
\label{thm:category-embedding}
Categories arise naturally from incidence structures closed under gluing operations.
\end{theorem}

\begin{proof}
\begin{itemize}
\item Objects are represented as incidences
\item Morphisms are represented as incidences with appropriate boundary structure
\item Composition is given by the gluing operation, which satisfies associativity (A8)
\item Identity morphisms are represented by unit incidences (A7)
\end{itemize}

The associativity of composition follows directly from \texttt{gluing\_associative} in \texttt{IncidenceTheory/Basic.lean}.
\end{proof}

\begin{theorem}[Type Embedding]
\label{thm:type-embedding}
Inductive and coinductive types can be represented through recursive incidence structures.
\end{theorem}

\begin{proof}
\begin{itemize}
\item Inductive types are represented by incidences that may refer to themselves through boundaries
\item Constructors are represented as incidences with appropriate argument structure
\item Coinductive types allow infinite unfoldings controlled by observational equivalence
\end{itemize}

This is formalized through \texttt{is\_inductive\_form}, \texttt{is\_coinductive\_form}, and related theorems in \texttt{IncidenceTheory/Basic.lean}.
\end{proof}

\section{Semantic Models}

We provide three semantic interpretations of the Theory of Incidence, establishing its relative consistency.

\begin{theorem}[ZF Set-Theoretic Model]
\label{thm:zf-model}
The Theory of Incidence can be consistently interpreted in ZF set theory.
\end{theorem}

\begin{proof}
Incidences are interpreted as sets with boundary functions:
\begin{itemize}
\item Each incidence $i$ is a set containing its boundary information
\item The boundary operator maps to finite sets of endpoint data
\item Gluing corresponds to set union with appropriate type constraints
\end{itemize}

This model preserves all axioms (A1–A17) and ensures well-foundedness through the finite boundary condition.
\end{proof}

\begin{theorem}[Adhesive Category Model]
\label{thm:adhesive-model}
Gluing operations correspond to pushouts in adhesive categories.
\end{theorem}

\begin{proof}
The gluing operation satisfies the properties of pushouts in adhesive categories:
\begin{itemize}
\item Van Kampen squares exist for gluing operations
\item The gluing preserves the incidence structure
\item Type consistency ensures category-theoretic properties
\end{itemize}
\end{proof}

\begin{theorem}[Homotopy Type Theory Model]
\label{thm:hott-model}
Recursive incidences correspond to coinductive types in HoTT.
\end{theorem}

\begin{proof}
\begin{itemize}
\item Inductive incidences model inductive types with path constructors
\item Coinductive incidences model coinductive types with guarded recursion
\item Observational equivalence corresponds to propositional equality in HoTT
\end{itemize}
\end{proof}

\begin{theorem}[Type Consistency]
\label{thm:type-consistency-semantic}
Boundary incidences share the same type as their parent incidence.
\end{theorem}

\begin{proof}
This follows directly from axiom A2, formalized as the \texttt{type\_consistent} property of the Incidence structure and proved as \texttt{type\_consistency} theorem in \texttt{IncidenceTheory/Basic.lean}.
\end{proof}

\begin{theorem}[Observational Equivalence as Bisimulation]
\label{thm:observational-bisimulation}
The relation $\approx$ defined by type and boundary equality forms a bisimulation.
\end{theorem}

\begin{proof}
Observational equivalence is reflexive, symmetric, and transitive, as established by theorems \texttt{approx\_refl}, \texttt{approx\_symm}, and \texttt{approx\_trans} in \texttt{IncidenceTheory/Basic.lean}. The bisimulation property ensures that equivalent incidences cannot be distinguished by their observable behavior.
\end{proof}

\section{Linear Semantics}
- Boundary matrices and Laplacians
- Linking logic to spectral methods

\section{Examples and Applications}
- Graph structures
- Self-referential systems
- Applications in computation and medicine

\section{Related Work}
Incidence Theory builds upon and extends several established areas:

- **Higher Operads and ∞-Categories**: Our recursive boundary operator and gluing align with the structure of higher operads (Leinster) and ω-categories, providing a concrete realization for abstract ∞-dimensional structures.
- **Adhesive Categories and DPO Rewriting**: Gluing semantics directly correspond to pushouts in adhesive categories, enabling formal model transformation.
- **Coalgebra and Coinduction**: The coinductive mode and observational equivalence draw from bisimulation in coalgebraic systems, adapted to relational structures.
- **Homotopy Type Theory**: Inductive/coinductive families in Inc mirror HoTT's types, with incidence boundaries as paths.

Inc differs by unifying these into a single relational primitive, avoiding meta-level machinery for recursion or self-reference.

\section{Conclusion}
The Theory of Incidence establishes a relational foundation unifying Set, Cat, and Type.

\appendix
\section{Lean Theorems and Axioms}
List of formalized axioms and theorems in Lean4 (see incidence-theory/ for code).

\end{document}
